% !TEX root =  master.tex
\chapter{KI‑gestützte nachhaltige Materialentwicklung}
...

\section{Toolsübersicht - Citrine}
Citrine Informatics ist ein auf Materialinformatik spezialisiertes Unternehmen, das Künstliche Intelligenz zur beschleunigten Entwicklung und Optimierung von Materialien einsetzt. Der Ansatz von Citrine basiert auf der Kombination physikbasierter Domänenkenntnisse mit maschinellen Lernverfahren, um aus heterogenen Materialdaten robuste Vorhersagemodelle abzuleiten. Ziel ist es, Zusammenhänge zwischen Zusammensetzung, Prozessparametern, Mikrostruktur und Materialeigenschaften systematisch zu erfassen und für das gezielte Design neuer Werkstoffe nutzbar zu machen. Durch diese datengetriebene Methodik können Entwicklungszyklen verkürzt, experimentelle Aufwände reduziert und nachhaltigere Materiallösungen schneller identifiziert werden. Insbesondere im industriellen Umfeld ermöglicht Citrine Informatics eine wissensbasierte Entscheidungsunterstützung, indem Materialanforderungen mit Leistungs‑, Kosten‑ und Nachhaltigkeitszielen in Einklang gebracht werden. \cite{citrine_informatics_platform_2026}

\section{APIs}%(mit exomatter Absatz zusammen schreiben?!)
%Software wird in der Regel über eine Benutzeroberfläche direkt von Menschen verwendet. Zunehmend wird Software jedoch nicht mehr ausschließlich von Endnutzern genutzt, sondern auch von anderen Softwaresystemen. Dies erfordert eine spezielle Form der Schnittstelle, die als Application Programming Interface (API) bezeichnet wird.
%APIs stellen einen standardisierten Mechanismus zur Anbindung, Integration und Erweiterung von Softwaresystemen dar. Sie werden insbesondere beim Aufbau verteilter Softwaresysteme eingesetzt, deren Komponenten lose miteinander gekoppelt sind. Die in diesem Zusammenhang betrachteten APIs sind überwiegend Web‑APIs, die als Webservices realisiert sind und Ressourcen über gängige Webtechnologien bereitstellen. Typische Anwendungsbeispiele für APIs finden sich in mobilen Anwendungen, Cloud‑basierten Systemen, Webanwendungen sowie in intelligenten Endgeräten.
%Ein wesentliches Merkmal von APIs ist ihre einfache, klare und wiederverwendbare Struktur, die eine effiziente Integration unterschiedlicher Anwendungen ermöglicht. Im Gegensatz zu klassischen Benutzeroberflächen sind APIs für Endnutzer in der Regel nicht sichtbar und werden nicht direkt von Menschen bedient. Stattdessen arbeiten sie im Hintergrund und werden von anderen Softwaresystemen oder Anwendungen aufgerufen. APIs dienen somit primär der Maschine‑zu‑Maschine‑Kommunikation sowie der Kopplung und Integration mehrerer Softwaresysteme.
%Die direkten Nutzer von APIs sind in der Regel Softwareentwickler, die Anwendungen oder Lösungen auf Basis dieser Schnittstellen implementieren. Aus diesem Grund müssen APIs mit einem klaren Fokus auf die Bedürfnisse der Entwickler konzipiert werden, die sie in neue oder bestehende Anwendungen integrieren. Diese Perspektive verdeutlicht die Notwendigkeit eines spezifischen, entwicklerzentrierten Ansatzes bei der Gestaltung und Implementierung von APIs. \cite{gu2016deep}
\subsection{REST API}
%Für diese Arbeit wird REST API benutzt um die Plattformen gemeinsam zu verbinden.

%REST (Representational State Transfer) ist ein architektonischer Ansatz zur Gestaltung von Diensten, die plattform‑ und umgebungsübergreifend genutzt werden können, um Interoperabilität im Kontext des World Wide Web zu gewährleisten. Zu den zentralen Merkmalen dieser Architektur zählen insbesondere Zustandslosigkeit sowie die Eignung für den plattformunabhängigen Zugriff. REST hat sich als weit verbreiteter und standardisierter Ansatz zur Bereitstellung von Diensten über das Internet etabliert. REST‑basierte Application Programming Interfaces (APIs) sind zudem ein integraler Bestandteil moderner Microservice‑Architekturen. Darüber hinaus existieren Forschungsarbeiten, die die REST‑Architektur zur Unterstützung verteilter Systeme erweitern.
%RESTful APIs, häufig auch als Web‑APIs bezeichnet, bestehen aus sogenannten Endpunkten. Jeder Endpunkt repräsentiert eine konkret implementierte Funktionalität eines Geschäftsprozesses. Der Zugriff auf diese Schnittstellen erfolgt in der Regel über das Hypertext Transfer Protocol (HTTP) unter Verwendung standardisierter Methoden wie GET, POST, PUT und DELETE. Die Identifikation und Adressierung der Ressourcen erfolgt dabei über eine eindeutige URI (Uniform Resource Identifier).
%Eine zentrale Herausforderung bei der Nutzung von REST‑Schnittstellen bestand lange Zeit in der Festlegung eines einheitlichen Nachrichtenformats für Anfragen und Antworten. Anfangs wurden REST‑APIs häufig informell beschrieben. Mit der Zeit etablierte sich jedoch JSON (JavaScript Object Notation) als de‑facto‑Standard für den Datenaustausch. JSON basiert auf einem textuellen Format und ist sowohl für Menschen gut lesbar als auch für Maschinen einfach zu verarbeiten, unabhängig von Plattform und Netzwerkumgebung. \cite{masse2011rest}







\cite{akafuah2016evolution}
\cite{tisza2013advanced}



